Un módulo de alimentación carga un condensador con corriente constante \(I>0\). Se adopta el modelo
\[
C v'=I-Gv,\qquad v(0)=0,
\]
donde \(C>0\) es la capacidad y \(G>0\) la conductancia de fuga. \(v\) se mide en voltios, \(t\) en segundos, \(C\) en faradios, \(G\) en siemens e \(I\) en amperios.

Compare oralmente tres diseños: el original, otro con el doble de capacidad y otro con el doble de conductancia, manteniendo los demás valores. Prediga la pendiente inicial y la tensión límite de cada uno, y dibuje a mano las tres respuestas en los mismos ejes. Justifique qué curvas pueden cruzarse después del arranque. No es necesario localizar el cruce ni desarrollar todas las fórmulas.
